\section{Procesos de preparación}

La preparación de la muestra es un paso crucial para aislar el analito y ponerlo en una forma adecuada para el análisis.

\subsection{Solubilización y digestión}
\begin{itemize}
    \item \textbf{Solubilización:} Disolver la muestra en un solvente adecuado.
    \item \textbf{Digestión:} Descomposición de la materia orgánica mediante ácidos o calor (ej. digestión con HNO$_3$ para metales).
\end{itemize}

\subsection{Procesos mecánicos}
\begin{itemize}
    \item \textbf{Molienda:} Reducir el tamaño de partícula para homogeneizar y facilitar la extracción.
    \item \textbf{Tamizado:} Separar la muestra por tamaño de partícula para obtener una fracción homogénea.
\end{itemize}

\subsection{Procesos de separación}
\begin{itemize}
    \item \textbf{Filtración:} Separar partículas sólidas de un líquido.
    \item \textbf{Extracción líquido-líquido (LLE):} Transferir el analito de una fase líquida a otra inmiscible.
    \item \textbf{Extracción sólido-líquido (SLE):} Extraer el analito de una matriz sólida con un solvente.
    \item \textbf{Extracción Soxhlet:} Técnica de extracción continua de sólidos con un solvente en reflujo.
\end{itemize}

\subsection{Dilución y concentración}
\begin{itemize}
    \item \textbf{Dilución:} Reducir la concentración del analito para que se encuentre dentro del rango de detección del instrumento.
    \item \textbf{Concentración:} Aumentar la concentración del analito para que sea detectable (ej. mediante evaporación del solvente).
\end{itemize}

\resumebox{ejercicio}{Seleccionar el pretratamiento}{\footnotesize
    Describe el pretratamiento necesario para determinar el contenido de cafeína en una bebida energética y para determinar el contenido de plomo en una muestra de pescado.
}

\resumebox{dato}{Extracción de compuestos bioactivos}{\footnotesize
    En la industria de alimentos y farmacéutica, se utilizan técnicas como la extracción Soxhlet o la extracción con fluidos supercríticos (CO$_2$) para obtener aceites esenciales o compuestos con propiedades antioxidantes de plantas.
}